% Source: Wald GR, physical PDF p. 212
%         (printed p. 199).
% Coverage: The neighborhood used in the proof of the strong-causality corollary.
% Figures/tables/footnotes: Figure 8.9; no tables or footnotes.
% Status: source-reviewed and compile-clean.
% Uncertainties: none.

\begingroup
\renewcommand{\thefigure}{8.9}
\begin{figure}[H]
  \centering
  \begin{tikzpicture}[x=.84cm,y=.84cm,>=Latex,line cap=round,font=\small]
    \draw[thick] (-2.55,-.25)
      .. controls (-2.75,.98) and (-2.00,1.78) .. (-1.12,1.66)
      .. controls (-.48,1.56) and (-.45,1.06) .. (-.17,1.25)
      .. controls (.02,1.37) and (-.08,2.20) .. (.57,2.35)
      .. controls (1.11,2.47) and (1.16,1.48) .. (1.50,1.30)
      .. controls (2.21,.96) and (2.15,.13) .. (1.84,-.62)
      .. controls (1.45,-1.56) and (.36,-1.86) .. (-.70,-1.58)
      .. controls (-1.78,-1.33) and (-2.46,-.87) .. (-2.55,-.25);
    \node at (1.15,.54) {\(V\)};
  \end{tikzpicture}
  \caption{The neighborhood, \(V\), of \(p\) used in the proof of the corollary to
  theorem~\ref{thm:8.2.2}.}
  \label{fig:8.9}
\end{figure}
\endgroup
